% Chapter 2, Literature Review.
% Generated from Dissertation/Chapter2_Literature_Review.md by tools/md2tex.py.
%% The draft-status line, the provenance notes and the [edit] markers in the
%% Markdown source are not part of the submitted text and are absent here.
%% Markdown is the source. Regenerate this file when it changes; do not hand-edit.

\chapter{Literature Review}
\label{ch:litreview}
\lhead{Chapter 2. \emph{Literature Review}}

The pharmacy clinical-trial set-up service accepts work as it arrives without explicit regard to available capacity, and some trials consequently miss the 60-day target. A planning response must identify feasible on-time commitments and limit lateness across the remaining queue. These are counted in different units, and the trust has not defined how to trade one against the other. Chapter 2 reviews the evidence relevant to both concerns and identifies the gap where they must be addressed together.

However, before work can be prioritised, the set-up process itself must be measured, and the literature provides little support for this. \citet{Walker2018Metrics} identify 117 trial-site performance metrics, none covering set-up or activation, and the eight core metrics \citet{Whitham2018Metrics} settle on by Delphi consensus all apply after a site opens. No established method therefore exists for measuring pharmacy set-up work, and this study must define the process and its measures before modelling it.

\section{Establishing an unrecorded process}
\label{sec:mapping}

Process mapping is the method for that step. \citet{Ofosu2021ProcessMapping} build a map of an emergency referral system from practitioner interviews and then examine it for gaps between the intended process and the actual one. This study follows the same procedure, and stops there because the trust does not record activity-level timings. Its output is an operational flowchart of how the three pharmacy teams and Research and Innovation interact.

The resulting map identifies a final approval step shared across trials, which makes commitment selection consequential where that resource is constrained.

\section{Selection under shared capacity}
\label{sec:selection}

Which trials to start is a recognised scheduling problem. Scheduling with rejection decides which jobs to accept as well as when to schedule them, because in a heavily loaded system accepting every job increases completion delays and tardiness costs (\citealp{Shabtay2013Rejection}).

\citet{SlotnickMorton2007OrderAcceptance} give the formulation this study follows. Their model selects jobs from a pool and balances the value of accepted jobs against their weighted tardiness, and early completion is neither penalised nor rewarded, so late completion carries a cost. \citet{LewisSlotnick2002MultiPeriod} extend the decision across planning periods, allowing acceptance decisions to be updated as new jobs arrive.

Similar models have been used in clinical research, but not for pharmacy trial set-up. \citet{DegliseHawkinson2020CAPTAIN} select trials for a clinical research unit that shares nursing and facility capacity. They showed that accepting an additional trial increases waiting times for other trials, and that this effect is non-linear, which is why they treat accepting a proposed trial as a decision that requires a model.

None of these models orders its objectives. \citet{SlotnickMorton2007OrderAcceptance} and \citet{LewisSlotnick2002MultiPeriod} combine theirs into one financial value, because revenue and lateness penalties are both money; \citet{DegliseHawkinson2020CAPTAIN} instead bound access delay and overtime. Each needs a number the trust cannot supply, a trade-off rate in the first case and a bound in the second.

\section{Ordered objectives}
\label{sec:ordered}

This study uses pre-emptive goal programming to order its objectives. In goal programming, each objective has a target, and the model seeks to reduce the difference between the target and the result. In the pre-emptive form, objectives are placed in order of importance. A lower-priority objective is considered only if it does not reduce the result achieved for a higher-priority objective (\citealp{JonesTamiz2010GoalProgramming}).

\citet{Munoz2018ProposalSelection} show when weighted goal programming is suitable: where decision-makers can state how important one objective is against another, as their clinical institute did through expert ratings. The trust has no such basis. \citet{BlakeCarter2002GP} instead used the pre-emptive approach in hospital planning, fixing an economically feasible case mix before choosing among the equally feasible alternatives. Section 3.5 applies that logic.

The model in this study has two distinctive features. First, its highest-priority objective is to maximise the number of trials completed on time, which is the measure the trust is assessed on. \citet{AtsmonyMosheiov2023OnTime} also maximise a weighted count of jobs meeting their due dates, although theirs is a just-in-time setting in which early completion is penalised too, and they use this as their only objective. Second, the model represents a single approval step that every trial must cross, whereas the machines in \citet{AtsmonyMosheiov2023OnTime} are independent. The next section explains why this step limits the number of trials that can be in progress.

\section{Modelling the shared step under uncertainty}
\label{sec:shared}

The approval step is a queue: trials wait, and one person completes the approval. \citet{Gupta2013Queueing} recommends the M/G/1 model for a single server whose service times follow no particular distribution, since its formula needs only the average service time and its variation, which are the only data the trust can provide. Appendix 4.B examines the arrival assumption it makes.

Similar queues have been used to model professional sign-off. \citet{LahiriSeidmann2012HangOvers} model a radiologist's workstation as a single-server queue, and \citet{Svirsko2022Medication} include waiting time that increases with workload in an optimisation model. Neither makes queue waiting time depend on which work is accepted: the first measures the queue without deciding which cases enter it, and in the second workload results from routing rather than selection decisions.

Section 3.4 brings approval waiting time into the trial-selection decision. The wait is shared by every committed trial and grows with the number committed, so a trial's on-time status depends on how many trials the plan commits, not on that trial alone.

The alternative was to simulate the full set-up pathway. Discrete-event simulation is widely used in NHS capacity planning (\citealp{Demir2017DemandCapacity,Harper2023Orthopaedic}), and \citet{Martinez2016TrialActivation} used it to model trial activation. The available records give coarse, retrospective segment timings, not the event-level activity, sequencing and service-time data such a simulation needs; they are sufficient for a queueing formula and not for a simulation of the full pathway.

One issue remains. The duration data are limited and highly variable, and on-time completion is a threshold outcome, so evaluating a plan on average durations does not show how often it will succeed. The selected plan is therefore tested through simulation, as \citet{Kheiri2021Theatre} test optimised NHS theatre schedules against bed-capacity breaches and \citet{Mavrotas2015Robustness} do for multi-objective models.

Two sources of uncertainty affect this test: the number of simulation runs, which sets the precision of the result (\citealp{Morris2019SimulationStudies}), and the duration distribution estimated from a small sample, which \citet{Barton2012InputUncertainty} calls input uncertainty. Chapter 4 reports both.

\section{Contribution of this dissertation}
\label{sec:contribution}

Table~\ref{tab:litsummary} sets the reviewed models against the problem's features.

\begin{table}[htbp]
  \centering
  \caption{Literature review summary}
  \label{tab:litsummary}
  \begin{tabularx}{\textwidth}{@{}lLLLLLLL@{}}
    \toprule
    \textbf{} & \textbf{Select} & \textbf{Shared step modelled} & \textbf{Value objective} & \textbf{Lateness} & \textbf{On-time count} & \textbf{Objectives} & \textbf{Setting} \\
    \midrule
    \citet{SlotnickMorton2007OrderAcceptance} & $\checkmark$ &  & $\checkmark$ & $\checkmark$ &  & Weighted & Manufacturing \\
    \addlinespace
    \citet{LewisSlotnick2002MultiPeriod} & $\checkmark$ &  & $\checkmark$ & $\checkmark$ &  & Weighted & Manufacturing \\
    \addlinespace
    \citet{BlakeCarter2002GP} & $\checkmark$ &  & $\checkmark$ &  &  & Pre-emptive & Hospital \\
    \addlinespace
    \citet{Munoz2018ProposalSelection} & $\checkmark$ &  & $\checkmark$ &  &  & Weighted & Clinical research institute \\
    \addlinespace
    \citet{Qi2022EquitableEnrolment} & $\checkmark$ &  & $\checkmark$ &  &  & Weighted & Clinical trials \\
    \addlinespace
    \citet{DegliseHawkinson2020CAPTAIN} & $\checkmark$ & $\checkmark$ & $\checkmark$ &  &  & Bounded & Clinical research unit \\
    \addlinespace
    \citet{Kheiri2021Theatre} &  &  &  &  &  & Single & NHS theatres \\
    \addlinespace
    \citet{Svirsko2022Medication} &  & $\checkmark$ &  &  &  & Frontier & Hospital pharmacy \\
    \addlinespace
    \citet{AtsmonyMosheiov2023OnTime} & $\checkmark$ &  &  &  & $\checkmark$ & Single & Parallel machines \\
    \addlinespace
    This study & $\checkmark$ & $\checkmark$ &  & $\checkmark$ & $\checkmark$ & Pre-emptive & NHS pharmacy \\
    \bottomrule
  \end{tabularx}
\end{table}

\noindent\footnotesize \emph{Note.} The value objective, lateness and on-time count columns record what each model optimises, not what it constrains.\normalsize

Two gaps follow. First, no reviewed model combines a shared step with pre-emptive objectives: \citet{DegliseHawkinson2020CAPTAIN} model a shared step but bound their other objectives rather than ordering them, and \citet{BlakeCarter2002GP} order their objectives but model no shared step. This dissertation closes that gap, modelling the shared approval step as an analytical queue whose waiting time is calculated for each possible plan size, and ordering two goals pre-emptively: the on-time count first, matching the measure the trust is assessed on, and weighted lateness second. Second, no reviewed model combines an on-time count with a value objective. That gap is not addressed here, because no trial in this analysis carries a priority score (Section~\ref{sec:inpractice}), and it remains open pending score coverage linked to set-up outcomes.

---

